% Part 3: the research literature, explained
\section{What the research literature says}
\label{sec:literature}

This section explains the studies a designer of \hv{} needs to understand,
at the level of what the researchers actually did and found. Each subsection
ends with the consequences for our design, so that
Section~\ref{sec:design} can cite conclusions rather than re-argue them. A
note on inspection depth: the load-bearing papers below were read at method
level (design, data, and headline numbers verified from the full text or
extended abstract); the supporting papers were read at abstract level and are
cited for orientation only. One result, StoryScope, reached us secondhand
through an open-source project and is marked accordingly wherever it appears.

\subsection{The fingerprint: what makes text read as machine-written}
\label{sec:lit-fingerprint}

Start with the study that made the phenomenon measurable.
\citet{kobak2025delving} asked how much of the biomedical literature is now
LLM-assisted, without using any detector. They parsed more than fifteen
million English PubMed abstracts from 2010--2024 and tracked, for every word,
the fraction of abstracts containing it each year,
\begin{equation}
p_{t}(w) \;=\; \frac{a_{t}(w)+1}{b_{t}+1},
\label{eq:kobak-p}
\end{equation}
where $a_t(w)$ counts abstracts in year $t$ containing word $w$ and $b_t$ is
the year's abstract total. Borrowing the excess-mortality design from
epidemiology, they project a counterfactual 2024 frequency from the pre-LLM
trend,
\begin{equation}
q_{2024}(w) \;=\; p_{2022}(w) \;+\; 2\,\max\bigl(p_{2022}(w)-p_{2021}(w),\,0\bigr),
\label{eq:kobak-q}
\end{equation}
that is, the 2021$\to$2022 change extrapolated two years forward and clamped
so the counterfactual never falls (2023 is skipped as already contaminated).
A word is in excess if its gap $\delta = p - q$ or its ratio $r = p/q$
clears a threshold. The result that matters for us is \emph{which} words
were in excess. The Covid years produced excess content words (nouns about
the pandemic); 2024 produced excess \emph{style} words: ``delves'' at a
ratio near 28, ``underscores,'' ``showcasing,'' ``potential'' with a gap of
0.045. From the share of abstracts containing at least one excess style
word, they bound LLM assistance below by 13.5\% of 2024 abstracts, up to
40\% in some subcorpora. The bound is a floor, since authors who scrub the
telltale words become invisible, and that caveat cuts both ways for us:
scrubbing words defeats this measurement without improving the prose.

Why do models overuse exactly these words? \citet{juzek2025delve} isolated
the vocabulary by a three-way intersection: words whose PubMed frequency
jumped significantly between 2020 and 2024, that ChatGPT also overuses when
asked to write abstracts for 2020 papers (so the comparison is
like-for-like), yielding 21 focal words. They then tested explanations. The
words are not overrepresented in plausible training corpora, which weakens
the training-data story. The sharpest test compares Llama-2-7B Base with its
Chat variant, identical architecture, the Chat model adding only supervised
fine-tuning and RLHF: the Chat model is markedly less surprised by
focal-word-laden text. The overuse enters at the human-feedback stage. One
more of their findings matters for design: raters recruited to resemble the
RLHF workforce showed no overall preference \emph{for} the focal words, so
nothing in the reward signal pushes back against amplification once it
starts.

For practitioners, the catalogue of the fingerprint is Wikipedia's ``Signs
of AI writing,'' maintained by the editors who clean AI text out of Wikipedia
\citep{wikipedia2026signs}. It goes well beyond vocabulary: inflated
significance (``stands as a testament''), superficial \mbox{-ing} analysis
tails, rule-of-three padding, negative parallelism (``not just X, but Y''),
formulaic challenges-and-future sections, essay-like conclusions, title-case
headings, and chatbot residue. Two properties make it the right seed list
for diagnostics. It is maintained against live cleanup experience, so it
updates as the models change; and its preamble insists that no single sign
is diagnostic, since human writers produce all of them at some rate. That
warning is the germ of the false-positive discipline that the best
open-source tools in Section~\ref{sec:oss} implement and the worst ignore.

What about ``slop'' as a quality construct rather than an authorship one?
\citet{shaib2025slop} asked nineteen experts across NLP, writing,
journalism, linguistics, and philosophy to define it, coded the definitions,
and landed on three themes: information utility (density and relevance),
information quality (factuality, bias), and style quality (repetition and
templatedness, coherence, tone). Professional copy-editors then annotated
news and retrieval-augmented QA passages, first with a binary slop judgment,
then with word-level spans tagged by code. Two findings carry over to us.
Regressions showed \emph{relevance and density}, not style, best predict the
binary slop judgment, which is a formal version of our record's lesson that
substance failures hide under style failures. And even trained annotators
agreed only moderately (theme-level Krippendorff's $\alpha$ of 0.34--0.45),
which calibrates how much precision we should expect from any judge, human
or model, on this construct.

Finally, homogenization. \citet{padmakumar2024writing} hired 38 writers to
produce short argumentative essays in three within-subject arms: alone, with
base GPT-3 suggestions, and with feedback-tuned InstructGPT suggestions, with
keystroke logging attributing every character to writer or model. Essays
co-written with the instruction-tuned model were measurably more similar to
one another (higher pairwise Rouge-L and BERTScore within topics, fewer
distinct n-grams and key points), the base model showed no such effect, and
the decomposition traced the homogenization to the model-contributed text
specifically. The instruction tuning, not the assistance, is the
homogenizer. \citet{doshi2024generative} found the same shape in a
creative-writing field experiment, individual quality up, collective
diversity down, and \citet{guo2023curious} showed the spiral direction:
models trained recursively on synthetic text lose lexical, syntactic, and
semantic diversity generation over generation. One discourse-level result is
worth logging despite its secondhand status: the StoryScope study is
reported to detect machine narrative at 93\% F1 from discourse features
alone, with professional surface rewriting moving detection by less than two
points \citep{russell2026storyscope}. We have not inspected the paper; if it
holds, it independently confirms that the durable fingerprint is structural,
which is also what our failure record found when 815 sentence repairs left a
volume unaccepted.

\paragraph{Consequences for the design.} The fingerprint is distributional
and structural, not a word list. So \hv{} measures excess against a genre
reference corpus, Kobak-style, with
equations~\eqref{eq:kobak-p}--\eqref{eq:kobak-q} applied to pre-2022
economics text rather than PubMed; it treats word lists as seeds for
measurement, not rewrite targets; it measures homogenization \emph{across}
our documents, not only within one; and it expects style repair alone to
leave the deeper fingerprint intact, which is why the acceptance layer is
reader-based.

\subsection{Why the models write this way}
\label{sec:lit-training}

Four studies, read together, explain the register our agents default to, and
the explanation changes what interventions can work.

\citet{kirk2024understanding} trained the same base models through the
alignment stages, supervised fine-tuning, best-of-$n$ sampling, and
PPO-based RLHF, on summarization and instruction-following, then measured
performance (GPT-4-judged win rates, in and out of distribution) and output
diversity (distinct n-grams, embedding dispersion, and an entailment-based
measure, per input and across inputs). RLHF wins on performance, in and out
of distribution, and pays for it with a substantial drop in per-input
diversity. Notably, turning the KL penalty up, the knob that is supposed to
keep the tuned model close to the base model, lowered performance without
restoring diversity. The diversity loss is not a mis-set hyperparameter; it
is what the objective buys.

Why does preference optimization sharpen so hard? \citet{zhang2025verbalized}
supply the cleanest mechanism: typicality bias in the preference data
itself. Human raters, for reasons cognitive psychology has documented for
decades (mere exposure, processing fluency), systematically favor familiar,
schema-congruent text. Model the reward as true utility plus
$\alpha \log \pi_{\mathrm{ref}}(y\,|\,x)$, a bonus for being typical under
the reference model. On 6{,}874 preference pairs whose two responses are
\emph{equally correct}, the fitted $\alpha$ is about 0.57--0.65 and highly
significant: raters reward typicality even with correctness held fixed.
Optimizing that reward provably sharpens the policy toward
$\pi_{\mathrm{ref}}^{\,1+\alpha/\beta}$, collapsing onto the mode whenever
many answers tie on utility, which is precisely the situation in prose
style. Their intervention is a prompt change: ask for \emph{five responses
with their probabilities} rather than one response. Requesting a verbalized
distribution recovers 1.6--2.1$\times$ the diversity of direct prompting on
creative tasks and about two thirds of the base model's diversity across
post-training checkpoints, at no retraining cost.

Verbosity has its own paper trail. \citet{singhal2023length} decomposed
where RLHF's reward gains come from on three preference datasets, and the
answer is mostly length. Within fixed length bins, reward improves little
(on one dataset, within-bin gains are a few percent of the total); shifting
mass toward longer outputs accounts for 70--90\% of the improvement. Their
sharpest demonstration replaces the learned reward with a pure length score:
PPO against ``longer is better'' reproduces most of the measured win-rate
gain of PPO against the real reward model. Padding is not a stylistic
accident; it is what the training objective paid for.
\citet{sharma2024sycophancy} complete the picture on the agreeable side:
preference-trained assistants systematically flatter and hedge because
raters prefer agreement, the same mechanism that seeds throat-clearing and
deference into technical prose.

\paragraph{Consequences for the design.} These are training-time artifacts,
so prompt-level pleading (``write naturally'') fights the model's objective
and loses; what works is changing the sampling question (verbalized
sampling), editing after generation, and measuring. Compression deserves
first-class status among edit operations, since verbosity is a paid-for
bias. And any LLM used \emph{inside} \hv{} as an editor or judge carries
these same biases, which is why the judging protocol below is designed
around them rather than around trust.

\subsection{Detection: how it works and why it cannot be a gate}
\label{sec:lit-detection}

\hv{} needs to understand detectors for two reasons: their failure modes
rule out an entire class of tempting designs, and their internal statistics,
used gently, are legitimate population-level instruments.

The detection literature is one idea refined three times. Machine text sits
where the model expects text to sit; human text does not. GLTR
\citep{gehrmann2019gltr} visualized this directly, coloring each token by
its rank in the predicted distribution, and lifted human detection accuracy
from 54\% to 72\% just by showing the ranks. DetectGPT
\citep{mitchell2023detectgpt} made it a statistic: perturb the passage
(mask spans, let T5 refill them, about a hundred variants), score original
and variants under the suspected source model, and flag when the original
sits at a local likelihood peak, since sampled text lies at peaks and human
text does not. It works (0.95 AUROC on their news benchmark) but needs
roughly the right source model and a hundred model calls per document.
Fast-DetectGPT \citep{bao2024fastdetect} removed the rewriting: one forward
pass gives the conditional next-token distribution at every position, and
the statistic compares the log-probability of each actual token with the
mean and spread of alternatives sampled from that same conditional,
\begin{equation}
\text{curvature} \;=\; \frac{\log p(x) - \tilde{\mu}}{\tilde{\sigma}},
\label{eq:fdg}
\end{equation}
with machine tokens scoring near the top of their conditionals (curvature
around 3) and human word choices looking like ordinary draws (around 0). A
surrogate scorer makes it black-box, it is 340$\times$ faster than
DetectGPT, and it flags 80\% of ChatGPT passages at a 1\% false-positive
rate on benchmarks. Binoculars \citep{hans2024binoculars} normalizes
differently: the ratio of a passage's perplexity under one model to its
cross-perplexity between two sibling models, which cancels the ``weird
prompt'' confound (their capybara problem) and detects across generators
with one global threshold, over 90\% detection at 0.01\% FPR on their
benchmarks.

Now the failures, which are the design-relevant part.
\citet{liang2023detectors} ran seven widely used detectors on two human
corpora: essays by US eighth-graders and TOEFL practice essays by non-native
speakers. The US essays passed nearly clean; the TOEFL essays were flagged
as AI at an average false-positive rate of 61\%, with 98\% flagged by at
least one detector. The mechanism is exactly the detectors' core statistic:
limited lexical range produces low perplexity, and low perplexity reads as
machine. Two prompt experiments nailed it: asking ChatGPT to enrich the
TOEFL essays' word choice dropped their false-positive rate to 12\%, and
asking it to simplify the US essays raised theirs to 57\%. Polished, formal,
conventional prose, which is to say \emph{scholarly} prose, sits on the
wrong side of this statistic by construction. \citet{weberwulff2023testing}
tested fourteen academic and commercial tools and found none reliable, with
accuracy below 80\% and easy degradation; \citet{krishna2023paraphrasing}
showed that one round of automatic paraphrase defeats the
perturbation-family detectors wholesale. Detection is unreliable exactly
where we would use it and trivially evaded by anyone trying.

There is one statistically respectable use, and it is population-level.
\citet{liang2024monitoring} estimate the \emph{fraction} $\alpha$ of a
corpus that is substantially LLM-modified without classifying any document.
Represent each document by which adjectives it contains; estimate each
adjective's occurrence probability under a human reference corpus ($P$,
pre-ChatGPT reviews) and an AI reference corpus ($Q$, generated reviews);
then choose $\alpha$ to maximize the mixture likelihood
\begin{equation}
\hat{\alpha} \;=\; \arg\max_{\alpha} \sum_{i} \log\bigl((1-\alpha)\,P(x_i) + \alpha\,Q(x_i)\bigr).
\label{eq:mixture}
\end{equation}
On constructed mixtures with known $\alpha$ the estimation error is under
two percent, an order of magnitude better and about $10^{7}$ times cheaper
than aggregating per-document detectors. Applied in the wild: 10.6\% of
ICLR~2024 review text substantially modified (up from 1.6\%), 16.9\% at
EMNLP, no change at Nature-family journals, and higher $\alpha$ in
late, low-confidence reviews. The estimator is meaningless for accusing an
individual document, by construction, which is exactly the property we
want.

\paragraph{Consequences for the design.} No per-document authorship verdict,
ever: the false-positive profile lands on precisely the writers and the
register we care about, and evasion is cheap anyway. What survives is
(i) curvature or Binoculars-style scores as \emph{corpus-level} smoothness
trends, advisory only, computed on detexed prose; and (ii) the mixture
estimator of equation~\eqref{eq:mixture} as the statistically honest way to
answer ``how much of our output still reads machine-typical this quarter,''
with $P$ estimated from pre-2022 economics text and $Q$ from our own
agents' raw drafts. Both enter \texttt{hv drift} in
Section~\ref{sec:proposal}, gated against per-document use.

\subsection{The edit paradigm: the strongest positive evidence}
\label{sec:lit-editing}

Most directly useful for \hv{} is
\citet{chakrabarty2024salvage}, because it tests the exact pipeline we
propose: generate, diagnose, edit. They generated about a thousand
paragraphs with three frontier models from instructions reverse-engineered
out of published prose, had professional editors mark them up, and
consolidated the editors' labels into seven idiosyncrasy categories:
clich\'e, unnecessary or redundant exposition, purple prose, poor sentence
structure, lack of specificity and detail, awkward word choice, and tense
inconsistency. Eighteen professional writers then produced 8{,}035
categorized span-level edits, the LAMP corpus. Three of their findings set
our expectations. Writing quality did not differ across the three models,
so the idiosyncrasies are an ecosystem property, not a vendor one. An
LLM-based span detector reached 0.46 precision against expert spans, but
the experts themselves only agreed at 0.57, so ``flag the bad spans'' is
inherently noisy for everyone and should feed a human or model editor
rather than an automatic rewrite. And the payoff experiment: in blind
rankings by the writers, edited AI text (whether the spans came from
experts or the detector) beat raw AI text decisively, raw output ranked
last 60\% of the time, while human-edited text still ranked first. Editing
recovers most of the gap; it does not close it.

Two smaller results refine the objective. \citet{cheng2025humt} built a
corpus-based metric of human-like tone and a steering method, and found
that users do not uniformly prefer maximal human-likeness: the optimum
depends on genre and task. ``As human as possible'' is the wrong target;
``appropriate to the register'' is the right one, which for us is fixed by
the field's own literature. And the style-transfer literature
\citep{jin2022style} is the cautionary map: pre-LLM transfer methods
chronically traded content preservation against style strength, which is
the trade our meaning-ledger and baseline invariants exist to refuse.

\paragraph{Consequences for the design.} Build the editor as
span-diagnose-then-edit with categories, not as one-shot rewriting; expect
span precision near 0.5 and design the workflow so a human or a second
model adjudicates flags; measure edit effect by blind pairwise comparison;
and never let an edit change a claim, which our existing meaning-ledger
skill already enforces in doctrine and \texttt{hv diff} will enforce in
code.

\subsection{Judging writing with models, without fooling ourselves}
\label{sec:lit-judging}

\hv{} needs cheap judgments of ``which version reads better,'' so the
LLM-as-judge literature is load-bearing, mostly for its list of traps.

\citet{zheng2023judging} established the baseline: on their multi-turn
benchmark, GPT-4's agreement with pooled expert votes was 85\%, against 81\%
human-human agreement, so a strong model is a usable proxy for preference
between chat answers. The same paper measured the biases operationally.
Position bias: judge two near-identical answers twice with order swapped;
every judge showed order-dependent flips. Verbosity bias: pad one answer
with rephrased duplicates of its own list items, adding no information;
two of three judges preferred the padded answer more than 90\% of the time.
Self-enhancement: models rate their own outputs higher, though their design
could not settle causality. \citet{panickssery2024llm} settled it with a
clean causal design on summarization: models recognize their own outputs
above chance out of the box, and fine-tuning a model to be \emph{better} at
self-recognition increases its self-preference linearly, while control
fine-tunes on unrelated tasks do not. Self-preference rides on
self-recognition, so the generator must never judge its own rewrite.
\citet{dubois2024length} showed the verbosity bias is largely removable by
regression: length-controlled win rates change model rankings and nearly
triple agreement with human evaluations on their benchmark. Finally,
\citet{chakrabarty2024art} calibrate optimism for prose specifically:
expert writers applying a creativity rubric failed LLM-written stories on
roughly ten times more rubric items than professional stories, while LLM
judges rated the same stories generously. For writing quality, model judges
drift lenient exactly where experts are harsh.

\paragraph{Consequences for the design.} The judging protocol writes
itself, negatively: pairwise, never absolute scores; order randomized or
both orders averaged; length-controlled; a different model family from the
generator and from the editor; one rubric dimension per call rather than
holistic ``is this good''; and periodic validation of the whole proxy
against recorded human verdicts, which our evaluation corpus contains. Any
result that survives those controls is worth acting on; nothing else is.

\subsection{Readability science: measure against the genre, not a formula}
\label{sec:lit-readability}

Classical readability formulas are the wrong instrument for this project,
and it is worth knowing why rather than just asserting it. Flesch-family
scores are functions of syllable counts and sentence lengths. LLM prose
often scores \emph{well} on them while failing readers, because its
failures live in cohesion, rhythm, and register, which the formulas do not
see. The modern evidence: the CLEAR corpus work collected thousands of
passages with human ease-of-reading judgments and showed learned models
beat the classical formulas substantially \citep{crossley2022clear};
Coh-Metrix demonstrated two decades ago that texts separate on cohesion and
discourse dimensions the formulas ignore \citep{graesser2004cohmetrix}; and
Biber's register work established that genres occupy different regions of a
multidimensional feature space, so ``natural'' is genre-relative
\citep{biber1988variation}. Two findings about scientific prose complete
the picture. Scientific abstracts have been getting harder to read for a
century, with jargon and sentence complexity rising steadily
\citep{plavensigray2017readability}, so the median published abstract is a
degrading target, not a gold standard. And jargon density carries a
professional price: titles and abstracts heavier in specialized terminology
collect fewer citations \citep{martinez2021terminology}, a result an
economist will recognize as an incentive argument for plain language, not
just an aesthetic one.

\paragraph{Consequences for the design.} Every \hv{} metric is referenced
to a target-genre human corpus (for us: pre-2022 economics working papers
and journal articles), never to an absolute scale; the diagnostic suite
spans lexis, sentence rhythm, and discourse cohesion rather than syllable
arithmetic; and ``write like the median abstract'' is explicitly not the
objective, which the policy's plain-language stance already reflects.

\subsection{Economics and publishing norms}
\label{sec:lit-econ}

Norms here are settled enough to build on. \citet{korinek2023generative},
the profession's de facto guide, treats LLM writing assistance as ordinary
productivity technology for economists, to be used and disclosed. Journal
policy has converged on disclosure rather than prohibition (Nature
portfolio: models cannot be authors, use must be documented; Science moved
from ban to disclosure). The prevalence work says assistance is already
everywhere: the mixture estimator of Section~\ref{sec:lit-detection} puts
substantial LLM modification at 7--17\% of recent ML conference review text
\citep{liang2024monitoring, liang2024mapping}; AI-assisted reviews
measurably shift scores and acceptance at ICLR \citep{latona2024lottery};
adoption in research writing skews toward non-native English speakers and
junior researchers, whose styles are converging \citep{lin2025divergent};
and undisclosed chatbot artifacts (``as an AI language model\ldots'')
appear verbatim in published papers \citep{glynn2024academai}.

\paragraph{Consequences for the design.} \hv{} is a quality project under
disclosure norms, full stop. Detector evasion is not a goal, would not be an
achievement (Section~\ref{sec:lit-detection} showed evasion is cheap), and
is ruled out as a success metric because detectors cannot define quality.
The prevalence work also hands us our reference corpora: pre-2022 economics
text is a clean human baseline; anything after 2023 is contaminated and
unusable for calibration.
